\section{The limiting quotient and a restricted positive construction}
\label{sec:geometry}

We first identify the limiting primitive under specified changes of
presentation. We then give a finite-scale positive construction with
additional hypotheses. Proposition~\ref{prop:negativelevi} rules out
extending its ambient positivity to all mixed presentations by the same
primitive.

\subsection{Relative primitives and elementary cancellations}
Fix Hermitian reference data on a surviving source $H$ and target $W$.
Let $C$ be an acyclic source sector and
$D=[0,F]:H\oplus C\to U$, where $F:C\to U$ is invertible.
Evaluate the primitive on metrics $h_V=h_H\oplus h_C$ and write
$A_j=[B_j,T_j]$. The metrics on $C,U$ and the map $F$ are fixed
throughout the static arrow-coordinate comparison.
Put $P_F=F^\dagger F$ and
$\lambda=\lambda_{\min}(P_F)>0$.
The zero-arrow acyclic primitive is
\begin{equation}\label{eq:acyclicprimitive}
 \mathcal S_{\ac}
 =z_0(\log\det h_C-\log\det h_U)+z_0\tr\Ein(P_F).
\end{equation}
With fixed determinant volume forms, subtracting this declared
expression from \eqref{eq:primitive} gives
\begin{equation}\label{eq:relativeprimitive}
 \begin{split}
 \mathcal S_{\mathrm{rel}}
 ={}&z_0\log\det h_H+z_1\log\det h_W-q_h(B)\\
 &-\sum_j\tr\bigl(T_jR(P_F)T_j^\dagger\bigr),\qquad
 q_h(B)=\sum_j\norm{B_j}_{\HS,h}^2 .
 \end{split}
\end{equation}
For fixed $F$ and acyclic metric data,
$0\leq R(P_F)\leq\lambda^{-1}\id$.
Thus the last term has absolute value at most
$\lambda^{-1}\sum_j\norm{T_j}_{\HS}^2$ and has
$O(\lambda^{-1})$ bounds for every fixed derivative in bounded
real $B,T$ coordinates. This is a relative static comparison, not
a constant subtraction on the full metric orbit.

Several elementary compatibilities follow from this exact formula.
For mutually orthogonal acyclic sectors with
$F=\diag_a(\Lambda_a\id)$, the error is
\begin{equation}\label{eq:sectorerror}
 -\sum_aR(\Lambda_a^2)\sum_j\norm{T_j^{(a)}}_{\HS}^2.
\end{equation}
The subtractions and removals of these sectors commute. Their bound
depends on $\min_a\Lambda_a$, not on the ratios of the scales.
This is an assertion for the displayed block decomposition, not
for interacting acyclic differentials without such a decomposition.

An elementary source shear
$g_V=\left(\begin{smallmatrix}\id&S\\0&\id\end{smallmatrix}\right)$,
where $S:C\to H$, leaves $D$ unchanged and replaces
$T_j$ by $T_j-B_jS$. The cohomology maps remain $B_j$.
For bounded data the two error terms in \eqref{eq:relativeprimitive}
both tend to zero, so their limits agree.
The same holds for every fixed derivative after composition with a
holomorphic parameter chart whose shear has bounded derivatives.
The two reference presentations must retain the same cohomology
Hermitian data, or transport them by a specified holomorphic isometry.
This does not assert equality at finite scale under a nonunitary shear
with unchanged ambient metrics.

Constant unitary frame changes preserve the comparison when they
transport $F$ and $T_j$ together. With unequal scales, a rotation of
acyclic frames generally changes the displayed matrix of $F$.
Rotating only the arrows while holding that matrix fixed is not the
claimed invariance.

Now specialize to $\dim H=\dim W=n>0$, $z_0=-1+\ii$, $z_1=1+\ii$.
The total cohomology charge is $2n\ii$, of phase $\pi/2$.
At fixed reference metrics, the limiting real arrow potential is
\begin{equation}\label{eq:limitpotential}
 K_H(B)=q(B)=\sum_j\norm{B_j}_{\HS}^2.
\end{equation}
Its Levi form is the usual positive flat form. The limiting
metric-direction function is
$q_h(B)+\log\det h_H-\log\det h_W$.
Its critical equations are
\begin{equation}\label{eq:ordinarybalance}
 \sum_jB_j^\dagger B_j=\id_H,\qquad
 \sum_jB_jB_j^\dagger=\id_W.
\end{equation}
The equality of dimensions is necessary by taking traces.
The resulting orbit-type K\"ahler quotients are the ordinary quiver
quotients. Their existence is classical; the calculation above
identifies the limiting potential and its normalization for these
presentations.

\subsection{An explicit finite-scale restriction}
Let $X=\operatorname{Hom}(H,W)^k$,
$G=\GL(H)\times\GL(W)$ and $K=\U(H)\times\U(W)$, acting by
$B_j\mapsto g_WB_jg_H^{-1}$. Take finitely many acyclic copies
$C_a=U_a=H$. Fix positive external scales $\Lambda_a$ and constant
matrices $L^{(a)}$ on the arrow-index space $\CC^k$ such that
\begin{equation}\label{eq:arrowmixing}
 (L^{(a)})^\dagger L^{(a)}=\kappa_a^2\id_k,\qquad \kappa_a\geq0.
\end{equation}
At least one $\kappa_a$ may be positive, so mixing need not vanish.
Define the holomorphic family of presentations
\begin{equation}\label{eq:graphfamily}
 T_j^{(a)}=\sum_lL_{jl}^{(a)}B_l,\qquad
 D=[0,\diag_a(\Lambda_a\id_H)].
\end{equation}
Its equivariance is for the embedding
\[
 g_V=\diag(g_H,g_H,\ldots,g_H),\qquad
 g_U=\diag(g_H,\ldots,g_H),\qquad g_W=g_W.
\]
It is not equivariant as a family with fixed graph data under the
entire presentation gauge group.

Transport $h_H$ to every $C_a$ and $U_a$, and retain $h_W$ on $W$.
All scales, arrow-index matrices and these fixed copy identifications
are held constant in the following tangent and metric comparisons.
Then $D^\dagger D$ on the $a$th copy is $\Lambda_a^2\id$ and
\[
 \sum_j\norm{T_j^{(a)}}_{\HS,h}^2
 =\kappa_a^2 q_h(B),\qquad
 \log\det h_V-\log\det h_U=\log\det h_H.
\]
Consequently the pullback of \eqref{eq:primitive}, after subtraction
of the constant $z_0n\sum_a\Ein(\Lambda_a^2)$, is exactly
\begin{equation}\label{eq:graphprimitive}
 \mathcal S_{\CC,\Lambda}
 =z_0\log\det h_H+z_1\log\det h_W-\alpha_\Lambda q_h(B),\qquad
 \alpha_\Lambda=1+\sum_a\kappa_a^2R(\Lambda_a^2)>0.
\end{equation}
This subtraction is constant on the specified tied-metric family,
as well as in $B$ at fixed scales. In contrast, it is not constant
on the full metric space of the complex.
The real arrow potential at phase $\pi/2$ is
$K_\Lambda=\alpha_\Lambda q$. Its metric and form are
$\alpha_\Lambda g_{\mathrm{flat}}$ and
$\alpha_\Lambda\omega_{\mathrm{flat}}$.
Their positivity follows directly from this formula.

\begin{proposition}\label{prop:quotientisometry}
For each fixed collection of data in
\eqref{eq:arrowmixing}--\eqref{eq:graphfamily}, the map
\begin{equation}\label{eq:scaling}
 \sigma_\Lambda:B\longmapsto\alpha_\Lambda^{-1/2}B
\end{equation}
induces an isometry on each connected smooth orbit-type stratum
from the ordinary balanced K\"ahler quotient to the quotient defined
by $K_\Lambda$ at the same central moment level.
It is the identity on the corresponding polystable cohomology
isomorphism classes.
The isometry preserves lengths, intrinsic distances and metric
completions on identically specified path domains.
\end{proposition}

\begin{proof}
Identify the Lie algebra dual with pairs of Hermitian matrices using
the fixed moment-map normalization
\[
 \mu_\alpha(B)=
 \left(-\alpha\sum_jB_j^\dagger B_j,\,
        \alpha\sum_jB_jB_j^\dagger\right),\qquad
 \eta=(-\id_H,\id_W).
\]
Let $Z=\mu_1^{-1}(\eta)$ and
$Z_\Lambda=\mu_{\alpha_\Lambda}^{-1}(\eta)$.
The diagonal scalar circle acts trivially; it may be retained as an
ineffective subgroup, and $\eta$ annihilates its Lie algebra since
the dimensions are equal. We do not assume freeness of the action.
These finite-dimensional holomorphic Hamiltonian reductions have
the usual smooth K\"ahler structures on connected orbit-type strata
\cite{sl91,sj95}. The quiver Kempf--Ness correspondence identifies
their points with polystable representations
\cite[\S6]{king94}. The hypotheses apply here: the ambient space is
flat K\"ahler, $K$ is compact, the level is central, and the source
Gram component controls $q(B)$, making the moment map proper.

The map $\sigma_\Lambda$ is a $K$-equivariant bijection from $Z$
to $Z_\Lambda$, preserving every stabilizer subgroup. It satisfies
\[
 \sigma_\Lambda^*(\alpha_\Lambda g_{\mathrm{flat}})
 =g_{\mathrm{flat}},\qquad
 \sigma_\Lambda^*(\alpha_\Lambda\omega_{\mathrm{flat}})
 =\omega_{\mathrm{flat}}.
\]
For clarity, if $v$ is a tangent lift and $\xi_X$ an infinitesimal
$K$-orbit vector, then
\[
 d\sigma_\Lambda(v)=\alpha_\Lambda^{-1/2}v,\qquad
 \xi_X(\sigma_\Lambda B)=\alpha_\Lambda^{-1/2}\xi_X(B).
\]
It follows directly that
\[
 \inf_{\xi\in\mathfrak k}
 \alpha_\Lambda\norm{d\sigma_\Lambda(v)+\xi_X(\sigma_\Lambda B)}^2
 =\inf_{\xi\in\mathfrak k}\norm{v+\xi_X(B)}^2.
\]
These are the squared reduced tangent norms. Stabilizer directions contribute
zero orbit vectors and cause no inverse estimate.
Furthermore $\sigma_\Lambda$ is the action of
$(\id_H,\alpha_\Lambda^{-1/2}\id_W)\in G$, so it leaves the
representation isomorphism class unchanged.

For the path statement, fix a $K$-invariant subset $\mathcal P\subset Z$
and use $\sigma_\Lambda(\mathcal P)$ on the other side.
Choose piecewise smooth paths within connected orbit-type strata,
with any allowed concatenations specified identically via
$\sigma_\Lambda$. The equality of tangent norms preserves every
admissible path length and hence their infimum on each finite-distance
component. If this convention first produces a pseudometric, divide
by zero distance on both sides. The resulting metric spaces, and
therefore their metric completions, are isometric.
No identification with a separate algebraic compactification is used.
\end{proof}

The finite-scale identity is elementary after \eqref{eq:graphprimitive};
it is not a general relative estimate for arbitrary mixing.
For the rank-two family \eqref{eq:balancedfamily}, scalar rescaling
preserves its stabilizer jump exactly, so this restricted quotient
identity remains valid without an orbit gap.

As a concrete example, take one copy and $L=\kappa\id_k$, $\kappa>0$.
Then $T_j=\kappa B_j$, and the off-diagonal block of the source
arrow Gram matrix is $E_H=\kappa C_H$.
It equals $\kappa\id$ on $Z$ and
$(\kappa/\alpha_\Lambda)\id$ on $Z_\Lambda$.
In either case it is nonzero, so the full source arrow Gram matrix
does not commute with $D^\dagger D=\diag(0,\Lambda^2\id)$.
This is mixing in the specified orthogonal Hermitian presentation.
The nonunitary shear $S=\kappa\id$ removes it algebraically, so
the example is not asserted to be an intrinsically nonsplit complex.

For a set $I$ of retained acyclic copies, put
$\alpha_I=1+\sum_{a\in I}\kappa_a^2R(\Lambda_a^2)$.
The calibrated map between the $I$ and $J$ levels is multiplication
by $\sqrt{\alpha_I/\alpha_J}$. These factors telescope under
successive removals, proving order independence for this construction.
Constant unitary changes of frames transport the differential and
arrows together as above.

\subsection{Limits of the comparison}
The comparison fixes each $\Lambda_a$ and each $L^{(a)}$.
Varying $\Lambda_a$ in a tangent direction differentiates
$\alpha_\Lambda^{-1/2}$ and introduces an additional term;
Proposition~\ref{prop:quotientisometry} makes no metric assertion in
such total-parameter directions.
It also does not claim that the full phase flow preserves the
tied-metric family. Indeed that family keeps the norm of $D$ fixed,
whereas the cancellation theorems prove its growth.

Theorem~\ref{thm:intromixed} applies to ordinarily balanced initial
tuples $B\in Z$ with any finite positive number $\ell$ of these acyclic
copies. Indeed the combined acyclic rank is $r=\ell n$, and
\[
 \sum_{a,j}\norm{T_j^{(a)}}_{\HS}^2
   =n\sum_a\kappa_a^2=:M^2,\qquad
 \sigma_{\min}(F_0)=\min_a\Lambda_a.
\]
If $\min_a\Lambda_a\geq\Lambda_*(n,M)$, the actual phase flow has
the theorem's cancellation and relaxation properties, independently
of the ratios of the initial scales. This applies for arbitrary $k$
and every stabilizer of the ordinarily balanced tuple.
The evolving metrics are those of the full phase flow, not metrics
constrained to remain in the family used for the quotient isometry.
The rescaled representative $\alpha_\Lambda^{-1/2}B$ is generally
not ordinarily balanced at the reference metrics, so the dynamical
theorem does not automatically apply to that representative.

Thus the general result of this section is the relative static
primitive limit under the declared elementary changes.
The finite positive metric is a restricted construction.
Neither result asserts a general metric on the moduli of all
complexes or an isometry under arbitrary derived equivalences.
